\section{Διαδραστικό Μυθιστόρημα (Interactive Fiction)}\label{sec:chapter2}

\subsection{Τι είναι το διαδραστικό μυθιστόρημα}
\subsubsection{Εισαγωγή}
\noindent
Το Interactive Fiction Technology Foundation ορίζει το διαδραστικό μυθιστόρημα (IF) ως είδος βιντεοπαιχνιδιού με το οποίο ο παίκτης αλληλεπιδρά κυρίως με τη χρήση κειμένου. Στο πλαίσιο αυτού του ορισμού εντάσσονται πολλαπλά έργα, με ενδιαφέρουσες και καινοτόμες μορφές ανά τις δεκαετίες. \\

\noindent
Τυπικά, τα διαδραστικά μυθιστορήματα είναι μη γραμμικά. Θα μπορούσαμε να παρομοιάσουμε την ιστορία ενός διαδραστικού παραμυθιού με ένα μονοπάτι. Σε συγκεκριμένα σημεία της ιστορίας, το μονοπάτι διακλαδώνεται, και ο παίκτης πρέπει να αποφασίσει ποια διαδρομή θα πάρει. Οι διαδρομές αυτές μπορούν να συνδέονται αργότερα στην πορεία της ιστορίας, ή να αποκλίνουν σε διαφορετικά μονοπάτια - λήξεις της ιστορίας. Με αυτόν τον τρόπο, δύο παίκτες είναι πιθανόν να βιώσουν ένα ξεχωριστό μονοπάτι από το ίδιο διαδραστικό μυθιστόρημα.\\

\noindent
Όπως αναφέρθηκε, τα διαδραστικά μυθιστορήματα βασίζονται κυρίως σε κείμενο για να δημιουργήσουν τον κόσμο τους και να καθοδηγήσουν τον παίκτη στην πλοκή του παιχνιδιού. Παρόλα αυτά, συχνά μπορεί να περιέχουν εικόνες, μουσική, animations, βίντεο και ηχητικά εφέ, για να “εμβυθίσουν” περισσότερο τον παίκτη στο σύμπαν τους.

\noindent
Τα παιχνίδια διαδραστικού μυθιστορήματος μπορούν να χωριστούν σε δύο μεγάλες κατηγορίες, με βάση τον μηχανισμό αλληλεπίδρασης με αυτά: Σε parser-based και choice-based. Σε ένα parser-based παιχνίδι οι παίκτες πληκτρολογούν εντολές φυσικής γλώσσας σε έναν προσομοιωμένο κόσμο και το παιχνίδι τις ερμηνεύει ως ενέργειες που πραγματοποιεί ο πρωταγωνιστής. Τα parser-based παιχνίδια συχνά περιέχουν παζλ που ο παίκτης πρέπει να λύσει για να προχωρήσει την ιστορία. Ως παζλ μπορεί να οριστεί οποιαδήποτε κατάσταση ή σενάριο στο οποίο ένα παραγωγικό επόμενο βήμα δεν είναι άμεσα εμφανές.\\

\noindent
Τα choise-based διαδραστικά μυθιστορήματα είναι μία ακόμα πολύ γνωστή μορφή διαδραστικού μυθιστορήματος. Σε αυτού του είδος διαδραστικών μυθιστορημάτων οι παίκτες αλληλεπιδρούν με το παιχνίδι μέσω επιλογών. Αυτές οι επιλογές μπορούν να παρουσιαστούν είτε ως μενού πολλαπλών επιλογών, είτε ως υπερσύνδεσμοι μέσα στο κείμενο της ιστορίας. Τα choise-based παιχνίδια τείνουν να εστιάζουν περισσότερο στην εξερεύνηση διακλαδούμενων αφηγήσεων και πολλαπλών καταλήξεων παρά σε γρίφους για επίλυση ή μυστικά για ανακάλυψη. Τα παιχνίδια που βασίζονται σε επιλογές μπορούν να είναι ψηφιακά, ή υλικά, όπως ένα βιβλίο τύπου Choose Your Own Adventure.\\

\subsection{Παιχνίδια Περιπέτειας}
\subsection{Ιστορική Αναδρομή}
\subsection{Πλατφόρμες Συγγραφής (Authoring Systems)}
\subsubsection{Τι είναι οι πλατφόρμες συγγραφής}
\subsubsection{Inform}
\subsubsection{TADS}
\subsubsection{Twine}
\subsubsection{ADRIFT}
\subsubsection{Άλλες πλατφόρμες συγγραφής}
\subsection{Διαγωνισμοί}
